\section{RAG Pipeline}

The RAG pipeline is built from several components, such as search, prompt builder, and others.
In this section, we examine each design choice and explain the rationale behind it.

\subsection{Search}
Text search plays a critical role in the RAG pipeline because retrieving relevant content and providing it to the LLM improves result consistency and reduces hallucinations.
In this thesis, we employ two types of search to achieve better results: vector search and lexical search.
Combining these two methods yields a hybrid approach that leverages the strengths of both.

\subsubsection{Vector Search}
After constructing the dataset, we need to retrieve relevant information for our RAG pipeline.
By default, MathHub offers a robust text-based search system; however, it does not support vector-based search.
Therefore, vector search functionality was necessary to enable semantic search over the content of MathHub.
We retrieved content from the URLs of MathHub and built a vector database for querying this content.

In this approach, we experimented with various embedding models to observe their impact on query results.
The results indicated that increasing the embedding dimensions beyond 1,000 can be counterproductive, yielding poor results.

\subsubsection{Lexical Search}
MathHub provides a text search function; however, this approach has a significant limitation. It performs pure text search,
meaning it only executes exact string matching rather than identifying keywords from a user query.
Furthermore, users tend to ask questions that describe a concept rather than using specific keywords.

For example, if the user input is ``What does an NP-hard problem mean?'' the pure text search will retrieve
unrelated content with low relevance scores because it requires the exact phrase ``NP-hard'' to appear in the results.

This limitation motivated us to develop a keyword extraction mechanism from the user input and apply the text search
on the extracted keywords.

\subsection{Gather Prerequisites}

Each item retrieved from the search is mapped to its specific prerequisites via the MathHub knowledge base. We fetch the
prerequisites and their content for inclusion in the prompt.

\subsection{Gather Symbols}